\section{XLA Evaluation}
\label{sec:evaluation}
\begin{figuretext}{figures/XLA rewrite case results.pdf}{XLA case coverage and attribute synthesis. The fragment classification follows concrete encodings; completeness guarantees additionally require the theorem's assumptions.}{fig:xla-results}{A downward tree partitions 209 compiled XLA rewrite cases into 119 semantically supported cases and 90 unsupported cases. Of the supported cases, 94 verify with explicit attributes, 76 require no RHS descriptors, and 43 require synthesis. Synthesis succeeds on 36 cases: all 10 classified within the theorem fragment and 26 of 33 beyond it. Seven cases remain unresolved.}
We evaluate attribute synthesis on 209 concrete rewrite cases representing
199 XLA rewrite schemas. Of these, 119 cases have the operation semantics
required by the verifier; the remaining 90 require semantics that are not
currently supported. Among the supported cases, 43 contain RHS attribute
descriptors to synthesize. The other 76 require no RHS descriptors: their
replacements use attribute-free operations, such as elementwise addition,
or directly return a matched value, as in
$\operatorname{add}(X,0) \to X$.
Figure~\ref{fig:xla-results} summarizes the results.
\end{figuretext}

With explicit RHS metadata formulas, the verifier proves 94 of the 119
supported cases (79.0\%). Attribute synthesis succeeds on 36 of the 43
eligible cases (83.7\%), inferring all RHS descriptors while preserving the
written rewrite conditions. These include all 10 cases classified within
the fragment of Theorem~\ref{thm:complete-synthesis} and 26 of 33 cases
beyond it. Within the fragment, completeness holds under the theorem's
stated assumptions. For the 26 beyond-fragment successes, independent
semantic verification checks the soundness of each synthesized rewrite
under the declared semantics and written conditions.

\begin{table}[!htbp]
\centering
\begin{minipage}{0.86\linewidth}
\caption{Specification size across 36 successfully synthesized cases representing 30 XLA rewrite schemas.}
\label{tab:xla-specification-size}
\begin{tabularx}{\linewidth}{@{}Xrr@{}}
\toprule
Representation & Total lines\textsuperscript{1} & Total tokens\textsuperscript{1}\\
\midrule
Corresponding XLA implementation excerpts & 1,398 & 10,849\\
TEPL with explicit RHS attributes & 243 & 3,370\\
TEPL with inferred RHS attributes & 198 & 2,277\\
\bottomrule
\end{tabularx}
\smallskip
\begin{minipage}{\linewidth}
\footnotesize
\textsuperscript{1} Counts exclude blank lines and comments.
\end{minipage}
\end{minipage}
\end{table}

Across these 36 successfully synthesized cases representing 30 XLA rewrite
schemas, attribute inference reduces TEPL specification size by 18.5\% in
lines and 32.4\% in tokens (Table~\ref{tab:xla-specification-size}). It
eliminates 47 explicit RHS descriptor occurrences while preserving the
rewrite equations and written conditions. The XLA code matches patterns,
transforms attributes, and builds replacement graphs. TEPL handles this
work through declarative rules and attribute inference, reducing the code
authors must write. Some XLA excerpts also cover extra cases and policies,
which contribute to their larger size.
